\scanpage{24}
\begin{center}6 - 42 -\end{center}
\hfill n° 392

2) Les hypothèses sur $A$ et $B$ étant les mêmes que dans l'exerc. 1, soit $C$ un anneau tel que $A\subset C\subset B$ et que le corps des fractions de $C$ soit égal à $L$ ; tout idéal $\neq0$ dans $C$ est alors un \uline{réseau} de $L$, (relativement à $A$).

a) Montrer que si $\uline{a}$, $\uline{b}$ sont deux idéaux de $C$ tels que $\uline{a}$ soit pseudo-inversible, on a $\chi(\uline{ab},C)=\chi(\uline{a},C)+\chi(\uline{b},C)$.

b) On dit que deux idéaux $\uline{a}$, $\uline{b}$ de $C$ sont \uline{pseudo-étrangers} si, pour tout idéal premier $\uline{p}$ de hauteur 1 dans $A$, $\uline{a}_{\uline{p}}$ et $\uline{b}_{\uline{p}}$ sont étrangers dans $C_{\uline{p}}$. Montrer que si $\uline{a}$ et $\uline{b}$ sont pseudo-étrangers et $\neq0$, on a $\chi(\uline{ab},C)=\chi(\uline{a},C)+\chi(\uline{b},C)$. (Se ramener au cas où $A$ est un anneau de valuation discrète, observer que l'on a $\chi(\uline{a},\uline{a}+\uline{b})=\chi(\uline{a}\cap\uline{b},\uline{b})$ et utiliser la relation (\quad) du n°\quad du § 4).

c) Déduire de b) que si $\uline{a}$ est un idéal $\neq0$ de $C$, et si $\uline{q}_i$, $(1\leq i\leq r)$, sont les idéaux primaires de la décomposition primaire réduite de $\uline{a}$, correspondant aux idéaux premiers non immergés, on a $\chi(\uline{a},C)=\sum_i\chi(\uline{q}_i,C)$.

d) Soient $\uline{q}$ un idéal primaire de $C$, $\uline{p}$ l'idéal premier correspondant. Si $\uline{p}$ n'est pas de hauteur 1, on a $\chi(\uline{q},C)=0$ ; dans le cas contraire, si $\uline{p}_0=\uline{p}\cap A$, on a $\chi(\uline{q},C)=\deg_A(\uline{q})\cdot\uline{p}_0$.

e) Soit $\uline{a}$ un idéal de $C$ ; pour que $\uline{a}$ soit pseudo-inversible, (exerc. 1 b)), il faut et il suffit que l'on ait $\chi(\uline{a},C)=\chi(\uline{a}B,B)$. (Même méthode que dans l'exerc. 1 b)).

3) Les hypothèses sur $A$, $B$ et $C$ sont les mêmes que dans l'exerc. 2. On dit encore que le réseau $C^*$ est la \uline{codifférente} de $C$, et que $C:C^*$, qui est un idéal de $C$, est la \uline{différente} de $C$ sur $A$ ; on la note $\uline{D}_{C/A}$.

a) Montrer que le \uline{conducteur} $\uline{f}_{C/B}$ de $B$ dans $C$, (chap. V, § 1, n°\quad), est égal à $B^*:C^*$, et que pour tout idéal premier $\uline{p}$ de hauteur 1 dans $A$, on a $(\uline{f}_{C/B})_{\uline{p}}C^*_{\uline{p}}=B^*_{\uline{p}}$, (utiliser le fait que $B_{\uline{p}}$ est un anneau principal).

b) Montrer que l'on a $\uline{D}_{C/A}\subset\uline{f}_{C/B}\cap\uline{D}_{B/A}$, (utiliser a)).

c) Montrer que, pour que $C^*$, considéré comme idéal fractionnaire de $C$, soit pseudo-inversible, (exerc. 1 b)), il faut et il suffit que pour tout idéal premier $\uline{p}$ de hauteur 1 dans $A$, on ait $B_{\uline{p}}(\uline{D}_{C/A})_{\uline{p}}=(\uline{f}_{C/B})_{\uline{p}}(\uline{D}_{B/A})_{\uline{p}}$, (utiliser a)).

d) Si $C^*$ est pseudo-inversible, tout idéal fractionnaire $\uline{a}\neq0$ de $C$ tel que $\uline{a}:\uline{a}=C$ est pseudo-inversible, (utiliser le fait que pour deux idéaux fractionnaires non nuls $\uline{a}$, $\uline{b}$ de $C$, on a $(\uline{ab})^*=\uline{a}^*:\uline{b}$).
\mspage{23}
